%=====================================================================
%  MASTER FILE — merged from main.tex + a2.tex + c2.tex
%  Utshav Saha (2305077)
%  Original section names / captions preserved verbatim.
%=====================================================================
\documentclass[11pt,a4paper]{article}

%---------------------------------------------------------------------
% PACKAGES (union of all three files + new additions)
%---------------------------------------------------------------------
\usepackage[margin=1in]{geometry}
\usepackage{graphicx}          % images
\usepackage{caption}           % \captionof outside floats
\usepackage{subcaption}        % side by side pic
\usepackage{wrapfig}           % wrapped figure
\usepackage{float}             % H float specifier
\usepackage{lipsum}            % for vugijugi lekha

\usepackage{amsmath, amssymb, amsthm}   % math + theorems
\usepackage{mathtools}         % \DeclarePairedDelimiter, \coloneqq
\usepackage{cancel}            % strike-out inside math
\usepackage{siunitx}           % units
\usepackage{bm}                % bold math

\usepackage[inline]{enumitem}  % nested/custom lists + inline lists
\usepackage{array, multicol, multirow}
\usepackage{booktabs}          % pro-looking rules
\usepackage{longtable}         % multi-page table
\usepackage[table]{xcolor}
\usepackage{colortbl}

\usepackage{listings}
\usepackage{algorithm}
\usepackage{algpseudocode}
\usepackage[normalem]{ulem}    % for strikethrough
\usepackage{tcolorbox}         % boxed notes
\usepackage{fancyhdr}          % headers/footers
\usepackage[numbers]{natbib}
\usepackage{hyperref}          % load late

%---------------------------------------------------------------------
% CUSTOM MARGIN ENVIRONMENT (from a2.tex)
%---------------------------------------------------------------------
\def\changemargin#1#2{\list{}{\rightmargin#2\leftmargin#1}\item[]}
\let\endchangemargin=\endlist

%---------------------------------------------------------------------
% COLORS + LISTINGS CONFIG (from main.tex)
%---------------------------------------------------------------------
\definecolor{codegreen}{rgb}{0,0.6,0}
\definecolor{codegray}{rgb}{0.5,0.5,0.5}
\definecolor{codepurple}{rgb}{0.58,0,0.82}
\definecolor{backcolour}{rgb}{0.95,0.95,0.92}

\lstset{
    backgroundcolor=\color{backcolour},
    commentstyle=\color{codegreen},
    keywordstyle=\color{magenta},
    numberstyle=\tiny\color{codegray},
    stringstyle=\color{codepurple},
    basicstyle=\ttfamily\footnotesize,
    breakatwhitespace=false,
    breaklines=true,
    captionpos=b,
    keepspaces=true,
    numbers=left,
    numbersep=5pt,
    showspaces=false,
    showstringspaces=false,
    showtabs=false,
    tabsize=2
}

%---------------------------------------------------------------------
% TABLE STYLING (from main.tex)
%---------------------------------------------------------------------
\setlength{\tabcolsep}{10pt}         % gap before text starts
\renewcommand{\arraystretch}{1.4}    % scale cell height
\setlength{\arrayrulewidth}{0.4mm}   % border width
\arrayrulecolor{blue}

%---------------------------------------------------------------------
% HYPERREF SETUP
%---------------------------------------------------------------------
\hypersetup{
    colorlinks = true,
    linkcolor  = blue,
    citecolor  = teal,
    urlcolor   = magenta
}

%---------------------------------------------------------------------
% CUSTOM COMMANDS / OPERATORS  (NEW)
%---------------------------------------------------------------------
\newcommand\kw[1]{\textbf{\itshape #1}}          % from main.tex
\newcommand{\R}{\mathbb{R}}                       % blackboard R
\newcommand{\Ex}[1]{\mathbb{E}\!\left[#1\right]}  % expectation
\newcommand{\dd}[2]{\frac{\partial #1}{\partial #2}} % partial deriv
\DeclareMathOperator*{\argmin}{arg\,min}          % new operator
\DeclareMathOperator{\sinc}{sinc}
\DeclarePairedDelimiter{\abs}{\lvert}{\rvert}     % \abs{x}, \abs*{x}
\DeclarePairedDelimiter{\norm}{\lVert}{\rVert}

%---------------------------------------------------------------------
% THEOREM ENVIRONMENTS  (NEW)
%---------------------------------------------------------------------
\theoremstyle{plain}
\newtheorem{theorem}{Theorem}[section]
\newtheorem{lemma}[theorem]{Lemma}
\theoremstyle{definition}
\newtheorem{definition}{Definition}[section]
\theoremstyle{remark}
\newtheorem*{remark}{Remark}

%---------------------------------------------------------------------
% HEADER / FOOTER  (NEW)
%---------------------------------------------------------------------
\pagestyle{fancy}
\fancyhf{}
\lhead{\LaTeX{} Master Compilation}
\rhead{Utshav Saha (2305077)}
\cfoot{\thepage}

%---------------------------------------------------------------------
\title{\textbf{\LaTeX} tutorial \\[0.4em]
       \large Mathematical Foundations of Image Synthesis \\
       \large CSE200: Online-1 on \LaTeX}
\author{Utshav Saha (2305077)}
\date{July 2026}

\begin{document}

\maketitle % title er jinispati dekhate

\begin{abstract}
\noindent This single document merges three earlier submissions: the \LaTeX{}
tutorial, the image-synthesis article, and the CSE200 online test. Beyond the
original content, it demonstrates additional machinery: theorem environments,
matrices, \texttt{booktabs} and \texttt{longtable}, \texttt{siunitx} units,
paired delimiters, custom operators, boxed notes, minipages, and multicolumn
text.
\end{abstract}

\tableofcontents
\listoffigures
\listoftables

\newpage

%=====================================================================
\part*{Part I --- \LaTeX{} Tutorial}
\addcontentsline{toc}{section}{Part I --- \LaTeX{} Tutorial}
%=====================================================================

\section{Introduction}
This is a
new document\footnote{with a footnote}

Some text with \emph{emphasis and \emph{nested} content}

Some text with \textit{italic and \textit{nested} content}

Some text with \textbf{italic and \textit{nested} content}

Something about \kw{apples} and \kw{oranges}.

\textcolor{red}{A colored text}
\sout{A strike through line}

\section{Title of the first section}

Text of material in the first section

Double newline for a new paragraph.

Third paragraph.

\subsection{Subsection of the first section}

Text of material in the subsection.

\section{Now we will learn about lists}

Lists are of two kinds - ordered and unordered
\vspace{2\baselineskip}

Ordered List
\begin{enumerate}
    \item Bus
    \item Truck
    \item Plane
\end{enumerate}
\vspace{2\baselineskip}

Unordered List
\begin{itemize}
    \item[!] Apple
    \item[-] Banana
    \item Orange
\end{itemize}

\vspace{2\baselineskip}

Nested List
\begin{enumerate}[label=(\roman*)]
  \item An entry
  \item Another One \begin{enumerate}
  \item this is at a
  \item this is at b
  \end{enumerate}
  \item Wow! Three entries
\end{enumerate}

\begin{description}
    \item[Dog:]  member of the genus Canis, which forms part of the wolf-like canids,
  and is the most widely abundant terrestrial carnivore.
\end{description}

\subsection{Extra list tricks (new)}
\texttt{enumitem} can compress spacing, restart counters, and set custom labels.

\begin{enumerate}[label=\textbf{Step \arabic*.}, leftmargin=*, itemsep=0pt]
  \item Load the scene description.
  \item Build the acceleration structure.
  \item Shade and write to the framebuffer.
\end{enumerate}

An inline list is also possible:
\begin{enumerate*}[label=(\alph*)]
  \item read,
  \item transform,
  \item write.
\end{enumerate*}

\newpage
\section{Playing with images}
\lipsum[1]

\vspace{2\baselineskip}
Figure types
\begin{itemize}
    \item \emph{\textbf{h}} ‘Here’ (if possible)
    \item \emph{\textbf{t}} Top of the page
    \item \emph{\textbf{b}} Bottom of the page
    \item \emph{\textbf{f}} A dedicated page only for floats
    \item \emph{\textbf{H}} Exactly here, no floating (needs \texttt{float})
\end{itemize}

\begin{figure}[ht]
    \centering
    % First Image
    \begin{subfigure}[b]{0.35\linewidth}
        \centering
        
\includegraphics[width=\linewidth]{cat3.jpg}
        \caption{hapi billu}
        \label{fig:cat3}
    \end{subfigure}
    \hspace{1 cm}
    % Second Image
    \begin{subfigure}[b]{0.35\linewidth}
        \centering
        
\includegraphics[width=\linewidth]{cat2.png}
        \caption{sad billu}
        \label{fig:cat2}
    \end{subfigure}

    \caption{Two states of Billu} % Optional: Main caption for both images
    \label{fig:both_cats}         % Optional: Main label
\end{figure}
\vspace{2\baselineskip}

In Figure~\ref{fig:cat2}, we can a sad billu and in Figure~\ref{fig:cat3}, we can a hapi billu and .
\\This is a new line.

\vspace{2\baselineskip}

\lipsum[2]

\subsection{Wrapped Figure}
\begin{wrapfigure}{r}{0.5\linewidth}
    \centering
    
\includegraphics[width=0.5\linewidth]{cat1.jpg}
    \caption{Caption}
    \label{fig:placeholder}
\end{wrapfigure}
\lipsum[2-3]

\subsection{Side-by-side minipages (new)}
Two \texttt{minipage}s let you place a figure and its explanation next to each
other without a \texttt{subfigure}.

\noindent
\begin{minipage}[t]{0.46\linewidth}
\vspace{0pt}
\centering

\includegraphics[width=\linewidth]{cat1.jpg}
\captionof{figure}{A minipage-hosted billu}
\label{fig:mini}
\end{minipage}
\hfill
\begin{minipage}[t]{0.48\linewidth}
\vspace{0pt}
The left minipage carries an image; this right one carries text. Both are boxes
of a fixed width, aligned by their top edges (\texttt{[t]}). This is the usual
way to build a two-column comparison inside an otherwise one-column document.
\end{minipage}

\section{Text size niye gutaguti}

{\tiny Ar vallagenaa porashuna korte :( }

{\small Ar vallagenaa porashuna korte :( }

{\large Ar vallagenaa porashuna korte :( }

{\Large Ar vallagenaa porashuna korte :( }

{\LARGE Ar vallagenaa porashuna korte :( }

{\huge Ar vallagenaa porashuna korte :( }

{\Huge Ar vallagenaa porashuna korte :( }

\subsection{Text Emphasis and Fonts}
This sentence contains \textbf{bold text}, \textit{italic text}, \emph{emphasized text}, and \underline{underlined text}. Font sizes can also vary such as {\Large Large text}, {\small small text}, and {\Huge huge text}.

\subsection{Font families and shapes (new)}
\textrm{Roman family}, \textsf{Sans-serif family}, \texttt{Typewriter family},
\textsc{Small Caps Shape}, \textsl{Slanted shape}, \uline{ulem underline},
\uwave{wavy underline}, \xout{crossed out}.

\section{Cholo Table banai (onek pera)}
\begin{tabular}{|l|m{3cm}|}
    \hline
   % This  & \parbox{3cm}{is rather long \\ new paragraph} \\ \hline
   This & is  \\ \hline
    a & table and it has a really really really really long text\\ \hline
\end{tabular}

\subsection{Aro complicated table}
\begin{tabular}{|c|c|c|}
\hline
    \multicolumn{3}{|c|}{This is a multi column} \\ \hline
    \rowcolor{yellow!50}1 & 2 &\cellcolor{black!10} 3 \\ \hline
    \multirow{2}{2cm}{A multi row} & 5 & 6 \\
     & 8 & 9 \\ \hline
\end{tabular}

\subsection{Booktabs table (new)}
Professional tables avoid vertical rules and use \verb|\toprule|,
\verb|\midrule|, \verb|\bottomrule|.

\begin{table}[H]
\centering
\arrayrulecolor{black}
\begin{tabular}{lrr}
\toprule
\textbf{Renderer} & \textbf{Frames/s} & \textbf{Memory (\si{\mega\byte})} \\
\midrule
Rasterizer      & 240 & 128 \\
Path tracer     &   4 & 512 \\
Hybrid          &  60 & 320 \\
\midrule
\textbf{Mean}   & \textbf{101.3} & \textbf{320.0} \\
\bottomrule
\end{tabular}
\arrayrulecolor{blue}
\caption{A clean \texttt{booktabs} table}
\label{tab:booktabs}
\end{table}

\subsection{Longtable, spanning pages (new)}
\begin{longtable}{|c|l|c|}
\hline
\textbf{ID} & \textbf{Primitive} & \textbf{Cost} \\ \hline
\endfirsthead
\hline
\textbf{ID} & \textbf{Primitive} & \textbf{Cost} \\ \hline
\endhead
\hline \multicolumn{3}{|r|}{\emph{continued on next page}} \\ \hline
\endfoot
\hline
\endlastfoot
1 & Sphere    & $O(1)$ \\ \hline
2 & Triangle  & $O(1)$ \\ \hline
3 & BVH node  & $O(\log n)$ \\ \hline
4 & Mesh      & $O(n)$ \\ \hline
5 & Volume    & $O(n^{3})$ \\ \hline
\caption{A \texttt{longtable} that can break across pages}
\label{tab:longtable}
\end{longtable}

\section{Cholo Math kori}
\begin{enumerate}
\item Let's begin with a formula $e^{i\pi} + 1 = 0$ But we can also do

$$ e = \lim_{n\to\infty} \left(1+\frac{1}{n}\right)^n = \lim_{n\to\infty}\frac{n}{\sqrt[n]{n!}}$$

\item We can do another:

$$ e = \sum_{n=0}^{\infty}\frac{1}{n!}$$

\item More Formulas:

$$  \vec{v} = <v_1,v_2,v_3>$$

$$\vec{v}\cdot \vec{w}$$

$$\begin{bmatrix}
    1 & 2 & 3 \\
    4 & 5 & 6
\end{bmatrix}$$

\item We want to write the solution for the quadratic equation $$\alpha x^2 + bx + c = 0$$.

The solution is
$$ x = \frac{1}{2\alpha} \times \left(-b + \sqrt{-b^2 \pm 4\cdot \alpha \cdot c}\right) $$

But there are two solutions actually.
\begin{align*}
 x_1 & = \frac{1}{2\alpha} \times \left(-b + \sqrt{-b^2 + 4\cdot \alpha \cdot c}\right)  \\
 x_2 & = \frac{1}{2\alpha} \times \left(-b - \sqrt{-b^2 - 4\cdot \alpha \cdot c}\right) \\
 & = \frac{-b - \sqrt{-b^2 - 4 \alpha c}}{2a}
\end{align*}
\end{enumerate}

\section{Hyperlinks}
Visit \href{https://www.google.com}{Google} for more information.
A bare URL: \url{https://www.overleaf.com/learn}.
An internal jump: see Section~\ref{sec:mathmod} or Table~\ref{tab:booktabs}.

\section{Pseudocode}
\begin{algorithm}[ht]
\caption{Calculate Factorial}
\label{alg:factorial}
\begin{algorithmic}[1] % The [1] adds line numbers to every line
    \Require $n \geq 0$
    \Ensure Returns $n!$

    \State $result \gets 1$

    \If{$n = 0$ \textbf{or} $n = 1$}
        \State \Return $result$
    \EndIf

    \While{$n > 1$}
        \State $result \gets result \times n$
        \State $n \gets n - 1$
    \EndWhile

    \State \Return $result$
\end{algorithmic}
\end{algorithm}

\begin{algorithm}[ht]
\caption{Ray--Scene Intersection (new example)}
\label{alg:raycast}
\begin{algorithmic}[1]
    \Require ray $r$, scene $S$
    \Ensure nearest hit $h$ or $\varnothing$
    \State $t_{\min} \gets \infty$; \quad $h \gets \varnothing$
    \ForAll{objects $o \in S$}
        \State $t \gets \Call{Intersect}{r, o}$
        \If{$t > 0$ \textbf{and} $t < t_{\min}$}
            \State $t_{\min} \gets t$
            \State $h \gets o$
        \EndIf
    \EndFor
    \State \Return $h$
\end{algorithmic}
\end{algorithm}

\section{Code Snippet}
\begin{lstlisting}[language=Python, caption=A simple Python function]
def calculate_fibonacci(n):
    """Calculate the nth Fibonacci number."""
    if n <= 0:
        return 0
    elif n == 1:
        return 1
    else:
        return calculate_fibonacci(n-1) + calculate_fibonacci(n-2)

print(calculate_fibonacci(10))
\end{lstlisting}

\noindent Inline code is also possible with \verb|\lstinline|: \lstinline[language=C]|for(int i=0;i<n;i++)|.

%=====================================================================
\clearpage
\part*{Part II --- Mathematical Foundations of Image Synthesis}
\addcontentsline{toc}{section}{Part II --- Mathematical Foundations of Image Synthesis}
%=====================================================================

\section{Introduction to Image Synthesis}
Image synthesis is the process of generating images from mathematical and algo-rithmic descriptions. Modern graphics systems rely on \textbf{precise mathematical models} and efficient computation. \emph{Even small numerical errors} can result in visible artifacts in rendered images. Typical formulations involve scalar values such as $x$, indexed parameters
like $I_{ij}$, and exponential terms $e^{x^2}$. These elements are combined to simulate realistic lighting behavior.

\section{Core Components}
\subsection{Rendering Pipeline}
The rendering pipeline consists of the following stages:

\begin{itemize}
    \item Scene Definition
    \item Camera Configuration
    \item[-] Light Placement
\end{itemize}

\begin{changemargin}{-1cm}{0cm}
\textbf{Important} Shading Model\end{changemargin}
\begin{itemize}
    \item[-] Diffuse Shading
    \item[-] Specular Highlights
\end{itemize}

\subsection{Execution Order}
\begin{enumerate}
    \item Input Processing
    \item Ray Construction
        \begin{enumerate}
            \item Primary Rays
            \item Secondary Rays
        \end{enumerate}
    \item Intersection Evaluation
        \begin{itemize}
            \item Object Hit Test
        \end{itemize}
\end{enumerate}

\section{Mathemetical Modeling}
Light interaction with surfaces can be expressed using the following equations.

\begin{equation}
    L = I_{ij} \cdot \cos(\theta) + R
\end{equation}

$$k = \frac{ab}{cd}$$

\begin{align*}
    f(x) &= x^2 + 2x + 1\\
        &= (x+1)^2\tag{2}
\end{align*}

\[
\lvert y \rvert =
\begin{cases}
    y & \text{if } x \geq 0 \\
    -y  & \text{if } x < 0
\end{cases}
\]

\[
p(t) =
\begin{cases}
    0 & \text{if } t < 0 \\
    e^-t  & \text{if } 0 \leq t < 3 \\
    ln(t-2) & \text{if } t \geq 3
\end{cases}
\]

$$ \mathcal{N}(x_i, \mu, \sigma) = \frac{1}{\sqrt{2\pi}\sigma}e^{-{\frac{ ({x_i-\mu})^2}{2\sigma^2}}} $$

\subsection{Matrices and vectors (new)}
Different matrix delimiters:
\[
A = \begin{pmatrix} a & b \\ c & d \end{pmatrix},\qquad
B = \begin{bmatrix} 1 & 0 \\ 0 & 1 \end{bmatrix},\qquad
C = \begin{vmatrix} a & b \\ c & d \end{vmatrix} = ad - bc,\qquad
D = \begin{Bmatrix} x \\ y \end{Bmatrix}.
\]

A general matrix with ellipses:
\[
M =
\begin{bmatrix}
m_{11} & m_{12} & \cdots & m_{1n} \\
m_{21} & m_{22} & \cdots & m_{2n} \\
\vdots & \vdots & \ddots & \vdots \\
m_{m1} & m_{m2} & \cdots & m_{mn}
\end{bmatrix}
\in \R^{m \times n}.
\]

The homogeneous perspective projection:
\[
\begin{bmatrix} x' \\ y' \\ z' \\ w' \end{bmatrix}
=
\begin{bmatrix}
f/a & 0 & 0 & 0 \\
0 & f & 0 & 0 \\
0 & 0 & \dfrac{z_f+z_n}{z_n-z_f} & \dfrac{2 z_f z_n}{z_n-z_f} \\
0 & 0 & -1 & 0
\end{bmatrix}
\begin{bmatrix} x \\ y \\ z \\ 1 \end{bmatrix}.
\]

\subsection{Multi-line and aligned equations (new)}
\verb|align| numbers every line; \verb|\notag| suppresses one.
\begin{align}
    L_o(\bm{x}, \bm{\omega}_o)
      &= L_e(\bm{x}, \bm{\omega}_o)
       + \int_{\Omega} f_r(\bm{x}, \bm{\omega}_i, \bm{\omega}_o)\,
         L_i(\bm{x}, \bm{\omega}_i)\,(\bm{\omega}_i \cdot \bm{n})\,
         \mathrm{d}\bm{\omega}_i \label{eq:rendering} \\
      &\approx L_e + \frac{1}{N}\sum_{k=1}^{N}
         \frac{f_r\, L_i\, (\bm{\omega}_k \cdot \bm{n})}{p(\bm{\omega}_k)} \notag
\end{align}

Equation~\eqref{eq:rendering} is the rendering equation. Use \verb|gather| when
lines need no alignment point:
\begin{gather}
    \nabla \cdot \bm{E} = \frac{\rho}{\varepsilon_0} \\
    \nabla \times \bm{B} = \mu_0 \bm{J} + \mu_0\varepsilon_0 \dd{\bm{E}}{t}
\end{gather}

Use \verb|split| for one long equation with a single number:
\begin{equation}
\begin{split}
    \Ex{X^2} &= \int_{-\infty}^{\infty} x^2 \, p(x)\,\mathrm{d}x \\
             &= \operatorname{Var}(X) + \left(\Ex{X}\right)^2.
\end{split}
\end{equation}

\subsection{Operators, delimiters, and cancellation (new)}
Custom operators: $\displaystyle \hat{\theta} = \argmin_{\theta \in \Theta} \norm{y - X\theta}_2^2$,
and $\sinc(x) = \dfrac{\sin \pi x}{\pi x}$.

Auto-sized delimiters: $\abs*{\dfrac{a}{b}}$ versus fixed $\abs{\frac{a}{b}}$.

Cancellation in a derivation:
\[
\frac{\cancel{2}\,a\,b}{\cancel{2}\,c} = \frac{ab}{c}.
\]

Units via \texttt{siunitx}: a wavelength of \SI{550}{\nano\metre}, an irradiance
of \SI{1361}{\watt\per\square\metre}, and an angle of \ang{45}.

\subsection{Theorems and definitions (new)}

\begin{definition}[Radiance]
\label{def:radiance}
Radiance $L$ is the radiant flux per unit projected area per unit solid angle,
measured in \si{\watt\per\square\metre\per\steradian}.
\end{definition}

\begin{theorem}[Energy conservation]
\label{thm:energy}
For any physically valid BRDF $f_r$ and all outgoing directions $\bm{\omega}_o$,
\[
\int_{\Omega} f_r(\bm{x},\bm{\omega}_i,\bm{\omega}_o)\,
(\bm{\omega}_i \cdot \bm{n}) \, \mathrm{d}\bm{\omega}_i \le 1 .
\]
\end{theorem}

\begin{proof}
The reflected flux cannot exceed the incident flux; dividing both sides by the
incident flux and applying Definition~\ref{def:radiance} gives the bound. Equality
holds only for a lossless surface.
\end{proof}

\begin{lemma}
A Lambertian BRDF $f_r = \rho/\pi$ satisfies Theorem~\ref{thm:energy} iff
$\rho \le 1$.
\end{lemma}

\begin{remark}
In practice $\rho \in [0.02, 0.9]$ for real-world dielectrics.
\end{remark}

\section{Performance Comparison}

\begin{table}[!h]
    \centering
    \begin{tabular}{|c|c|c|}
    \hline
         Method & Accuracy & Time Complexity  \\ \hline
         Rasterization & Medium & $O(n)$ \\ \hline
        \multirow{2}{2cm}{Ray Tracing} & High & $O(n^2)$ \\
        & Very High & $O(n^3)$ \\ \hline
        \multicolumn{2}{|c|}{Hybrid Method} & $O(n^2)$ \\ \hline
    \end{tabular}
    \caption{Comparison of Image Synthesis Techniques}
    \label{tab:comparison}
\end{table}

The results in Table \ref{tab:comparison} highlight the trade-off between accuracy and computational cost.

\begin{tcolorbox}[colback=blue!5!white, colframe=blue!60!black,
                  title=Takeaway (new: tcolorbox)]
Rasterization wins on throughput; ray tracing wins on correctness. Hybrid
deferred pipelines are the usual compromise in real-time engines.
\end{tcolorbox}

\clearpage

\section{Visual Illustration}
\begin{figure}[ht]
    \centering
    \begin{subfigure}[b]{0.3\linewidth}
        \centering
        \includegraphics[width=1\linewidth]{A (1).PNG}
        \caption{Output A}
        \label{fig:outA}
    \end{subfigure}
    \hfill
    \begin{subfigure}[b]{0.3\linewidth}
        \centering
        \includegraphics[width=1\linewidth]{B (1).PNG}
        \caption{Output B}
        \label{fig:outB}
    \end{subfigure}
    \hfill
    \begin{subfigure}[b]{0.3\linewidth}
        \centering
        \includegraphics[width=1\linewidth]{C (1).PNG}
        \caption{Output C}
        \label{fig:outC}
    \end{subfigure}

    \caption{Comparison of three rendering outputs.}
    \label{fig:outputs}
\end{figure}

Figure \ref{fig:outputs} compares three different rendering outputs arranged in a single row.

\section{Conclusion}
Image synthesis combines \underline{mathematical precision}, algorithmic design, and efficient implementation. As scene complexity increases, choosing the appropriate method becomes critical for balancing quality and performance.

%=====================================================================
\clearpage
\part*{Part III --- CSE200: Online-1 on \LaTeX}
\addcontentsline{toc}{section}{Part III --- CSE200: Online-1 on \LaTeX}
%=====================================================================

\section{Introduction to Document Preparation}

Document preparation systems are essential for producing structured and professional technical documents. \textbf{LATEX} allows authors to focus on content while maintaining consistent formatting. This article demonstrates formatted text, lists, equations, tables, and figures using commands taught in class.

\section{List Structures}
\subsection{Mixed and Nested List}
\begin{itemize}
    \item Primary Feature
    \item Secondary Features
        \begin{itemize}
    \item[-] Inline Math
    \item[-] Display Math
        \end{itemize}
\end{itemize}

\textbf{Note} Lists can be customized.

\section{Mathematical Modeling}
\label{sec:mathmod}
Mathematical expressions are commonly used to describe relationships among
variables. Let the variables be $x_i$, $y_j$, and $\sigma^2$. Consider also
$x_i, y, \text{ and } z^2.$

\subsection{Displayed Equations}

\begin{equation}
    y_j = \alpha{x_j}^2 + \beta x_j + \gamma
\end{equation}

\begin{align*}
     T &= \sum_{i=1}^{n}(x_i - \mu)^2 \\
     &= \sum_{i=1}^{n}x_i^2 - 2\mu\sum_{i=1}^{n}x_i + n\mu^2 \tag{2}
\end{align*}

\[
f(x_i) =
\begin{cases}
    \frac{{x_i}^2}{\sigma^2} & \text{if } x_i \geq 0 \\
    \frac{|x_i|}{\sigma}  & \text{if } x_i < 0
\end{cases}
\]

\subsection{Two-column math (new)}
\begin{multicols}{2}
\noindent
Left column:
\[
\int_0^1 x^n \, \mathrm{d}x = \frac{1}{n+1}
\]
\[
\sum_{k=0}^{n} \binom{n}{k} = 2^n
\]

\columnbreak

\noindent
Right column:
\[
\prod_{i=1}^{n} i = n!
\]
\[
\lim_{x \to 0} \frac{\sin x}{x} = 1
\]
\end{multicols}

\section{Tabular Data Presentation}

\begin{table}[!h]
    \centering
    \begin{tabular}{|c|c|c|}
    \hline
        \multirow{2}{2cm}{\centering \textbf{Module}} & \multicolumn{2}{|c|}{\textbf{Score}} \\ \cline{2-3}
          & Theory & Lab  \\ \hline
          Module A & 75 & 80 \\ \hline
          \multirow{2}{2cm}{\centering Module B} & 88 & 90 \\
          & 85 & 87 \\ \hline
          Module C & 70 & 72 \\ \hline
    \end{tabular}
    \caption{Module-wise Performance Summary}
    \label{tab:modules}
\end{table}

\subsection{Aligned-decimal table with siunitx (new)}
\begin{table}[H]
\centering
\arrayrulecolor{black}
\begin{tabular}{l S[table-format=3.2] S[table-format=2.1]}
\toprule
{\textbf{Module}} & {\textbf{Mean}} & {\textbf{Std.\ dev.}} \\
\midrule
Module A & 77.50 & 3.5 \\
Module B & 87.50 & 1.3 \\
Module C & 71.00 & 14.0 \\
\bottomrule
\end{tabular}
\arrayrulecolor{blue}
\caption{Decimal-aligned summary using \texttt{siunitx} S columns}
\label{tab:siunitx}
\end{table}

\clearpage

\section{Visual Comparison}
\subsection{Asymmetric Subfigure Layout}
    \begin{figure}[ht]
    \centering
    \begin{subfigure}[b]{0.8\linewidth}
        \centering
        \includegraphics[width=1\linewidth]{A (1).PNG}
        \caption{Main Diagram}
        \label{fig:mainDiag}
    \end{subfigure}
    \begin{subfigure}[b]{0.3\linewidth}
        \centering
        \includegraphics[width=1\linewidth]{B (1).PNG}
        \caption{Detail B}
        \label{fig:detailB}
    \end{subfigure}
    \hfill
    \begin{subfigure}[b]{0.3\linewidth}
        \centering
        \includegraphics[width=1\linewidth]{C (1).PNG}
        \caption{Detail C}
        \label{fig:detailC}
    \end{subfigure}

    \caption{Asymmetric layout with one dominant figure}
    \label{fig:asymmetric}
\end{figure}

Figure \ref{fig:asymmetric} illustrates how visual elements can be arranged systematically within a document.

\section*{Conclusion}
\addcontentsline{toc}{section}{Conclusion}
This article highlighted the importance of \textbf{structured writing} using \LaTeX.The use of \textit{formatted text}, mathematical expressions, tables, and figures improves both \underline{readability} and presentation quality.

%=====================================================================
\section{Citation Examples}

I am citing an example work, \cite{vaswani2017attention, 10.1093/bioinformatics/btu462}

% \nocite{}

\citet{vaswani2017attention} introduced transformers in \citeyear{vaswani2017attention}

\bibliographystyle{plainnat}
\bibliography{biblio}

\appendix
\section{Appendix: Symbol Table}
\begin{table}[H]
\centering
\arrayrulecolor{black}
\begin{tabular}{cl}
\toprule
\textbf{Symbol} & \textbf{Meaning} \\
\midrule
$L_o$        & Outgoing radiance \\
$f_r$        & Bidirectional reflectance distribution function \\
$\bm{\omega}$& Direction on the unit sphere \\
$\rho$       & Albedo \\
$\sigma$     & Standard deviation \\
\bottomrule
\end{tabular}
\arrayrulecolor{blue}
\caption{Notation used throughout this document}
\label{tab:symbols}
\end{table}

\end{document}